\sloppy

% Настройки стиля ГОСТ 7-32
% Для начала определяем, хотим мы или нет, чтобы рисунки и таблицы нумеровались в пределах раздела, или нам нужна сквозная нумерация.
\EqInChapter % формулы будут нумероваться в пределах раздела
\TableInChapter % таблицы будут нумероваться в пределах раздела
\PicInChapter % рисунки будут нумероваться в пределах раздела

% Добавляем гипертекстовое оглавление в PDF
\usepackage[
bookmarks=true, colorlinks=true, unicode=true,
urlcolor=black,linkcolor=black, anchorcolor=black,
citecolor=black, menucolor=black, filecolor=black,
]{hyperref}

% Изменение начертания шрифта --- после чего выглядит таймсоподобно.
% apt-get install scalable-cyrfonts-tex

\IfFileExists{cyrtimes.sty}
    {
        \usepackage{cyrtimespatched}
    }
    {
        % А если Times нету, то будет CM...
    }

\usepackage{graphicx}   % Пакет для включения рисунков
\usepackage{adjustbox}  % Таблицы не шире \linewidth
\usepackage{float}
\usepackage{pdfpages}
\usepackage{tabu}
\usepackage{tabto}
% С такими оно полями оно работает по-умолчанию:
% \RequirePackage[left=20mm,right=10mm,top=20mm,bottom=20mm,headsep=0pt]{geometry}
% Если вас тошнит от поля в 10мм --- увеличивайте до 20-ти, ну и про переплёт не забывайте:
\geometry{right=20mm}
\geometry{left=30mm}


% Пакет Tikz
\usepackage{tikz}
\usetikzlibrary{arrows,positioning,shadows}

% Произвольная нумерация списков.
\usepackage{enumerate}

% ячейки в несколько строчек
\usepackage{multirow}

% itemize внутри tabular
\usepackage{paralist,array}

\usepackage[linesnumbered,ruled,vlined]{algorithm2e}

\usepackage[obeyFinal]{todonotes}
\usepackage{lastpage}
 
\newcommand{\AKHworries}[1]{\todo[inline, color=red!20]{#1}}
\newcommand{\YEYUworries}[1]{\todo[inline, color=violet!20]{#1}}
\newcommand{\AGworries}[1]{\todo[inline, color=blue!20]{#1}}
\newcommand{\RSworries}[1]{\todo[inline, color=brown!20]{#1}}

\usepackage{setspace}
 \usepackage{pdflscape}

% Перенос длинных подписей к таблицам/рисункам (ГОСТ + узкие поля)
\usepackage{caption}
\captionsetup{
    width=0.92\linewidth,
    justification=raggedright,
    font=small,
    labelfont=bf,
    singlelinecheck=false,
}
% Подпись к таблице не шире полосы набора (избегает вылета вправо)
\captionsetup[table]{width=\linewidth}

% Заголовки ячеек в две строки при узких столбцах
\usepackage{makecell}
\renewcommand\theadfont{\bfseries\normalsize}